\documentclass[Markos_Delaportas_1316.tex]{subfiles}

\begin{document}

\section{Experimental Results and Discussion}
This section consolidates the empirical and simulated findings of the
thesis, structured
to provide a comparative analysis of the Over-the-Air (OTA)
computation performance across
different modeling layers. The discussion first examines the
degradation mechanics revealed
by MATLAB and Simulink simulations, evaluating the efficacy of
applied equalization techniques.
It then outlines the evaluation framework for the physical
Software-Defined Radio (SDR)
deployments.

\subsection{Comparative Analysis of Channel Impairments (MATLAB vs.\ Simulink)}
The rigorous simulation of \gls{ota} architectures across
discrete-time algorithmic environments
(MATLAB) and continuous-time dynamical environments (Simulink)
provided deep insights into
the physical limitations of signal superposition.

Under ideal \gls{awgn} conditions, both simulation environments
confirmed the theoretical Mean
Squared Error (MSE) scaling laws, demonstrating that computation
variance is strictly bounded
by the channel noise \cite{sahin_survey_2023}. However, the
introduction of uncoordinated
physical impairments revealed significant divergence from the idealized theory:
\begin{itemize}
    \item \textbf{Multipath and Fractional Delay:} The high-fidelity
        simulations demonstrated that even sub-symbol fractional
        delays (modeled via Farrow filters) induce severe Intersymbol
        Interference (ISI). The 2D and 3D MSE surface scans revealed
        an irreducible error floor that scales with the multipath
        gain, fundamentally bottlenecking the computation accuracy
        regardless of the transmission power.
    \item \textbf{Asynchronous Carrier Frequency Offset (CFO):}
        Independent local oscillator offsets among distributed
        transmitters were shown to convert stationary baseband values
        into beat-frequency harmonics. Simulink evaluations
        demonstrated that traditional coherent detection fails
        catastrophically under these conditions. The application of
        non-coherent Euclidean magnitude detection mitigated the
        zero-crossing signal collapse but introduced a multi-user
        vector fading penalty.
\end{itemize}

\subsection{Evaluation of Equalization and Recovery Strategies}
The simulations also served to benchmark classical equalization and
synchronization recovery
schemes. The divergence between algorithmic MATLAB models (which
    assume perfect block-level
knowledge) and Simulink models (which emulate real-time causal
processing) highlighted the
difficulty of correcting asynchronous CFO post-superposition. While
techniques such as
squaring loops and Costas loops are standard for point-to-point
links, their application to
aggregated \gls{ota} signals proved inadequate when phase offsets are
uncorrelated across
multiple transmitters. This finding justifies the ongoing
investigation into more robust,
adaptive filtering mechanisms.

\subsection{Framework for SDR Experimental Validation}
Building upon the simulation baseline, the physical deployment
utilizing GNU Radio and \gls{sdr}
hardware aims to validate the theoretical error floors in a
real-world electromagnetic environment.
The planned experimental results will be benchmarked against the
Simulink diagnostics to quantify
hardware-specific non-linearities, such as amplifier saturation and
discrete sampling jitter.
Ultimately, the experimental framework will demonstrate the practical
computation rate achieved
by the \gls{ota} approach relative to traditional orthogonal access
schemes (e.g., \gls{tdma}),
illustrating the critical balance between resource savings and analog
computation error.

\ifSubfilesClassLoaded{
    \clearpage
    \printbibliography
}{}

\end{document}
